\documentclass[10pt,a4paper]{amsart}
\usepackage[margin=19mm]{geometry}
\usepackage{amsmath,amssymb,mathtools}
\usepackage[scr=rsfs,cal=euler]{mathalfa}
\usepackage{xfrac}
\usepackage[unicode,hidelinks,bookmarks=false]{hyperref}
\newcommand{\ie}{i.e.\ }
\newcommand{\cf}{cf.\ }
\newcommand{\resp}{respectively\ }
\newcommand{\Subsection}{\subsection}
\providecommand{\nobreakdash}{\nobreak\hbox{-}}
\setlength{\emergencystretch}{2em}
\multlinegap0pt
\usepackage{mdframed}
\usepackage{mdframed}
\usepackage{mdframed}
\usepackage{mdframed}
\usepackage{amssymb}
\usepackage{mathtools}
\usepackage{tikz}
\usepackage{mdframed}
\usepackage{tcolorbox}
\usepackage{enumerate}
\usepackage[shortlabels]{enumitem}
\newtheorem{lemma}{Lemma}
\newcommand{\C}{\mathbb{C}}
\newcommand{\covol}{\on{covol}}
\newcommand{\mc}{\mathcal}
\newcommand{\on}{\operatorname}
\title{Remarks on the disproof of the unit distance conjecture}
\author{Noga Alon, Thomas F. Bloom, W. T. Gowers, Daniel Litt, Will Sawin, Arul Shankar, Jacob Tsimerman, Victor Wang, Melanie Matchett Wood}
\date{Source: arXiv:2605.20695v1 (CC BY 4.0)}
\begin{document}
\maketitle

\noindent\emph{Source and licence.} Remarks on the disproof of the unit distance conjecture. Version arXiv:2605.20695v1 (CC BY 4.0), distributed by arXiv under the Creative Commons Attribution 4.0 International licence, http://creativecommons.org/licenses/by/4.0/, as recorded in the arXiv licence metadata for this exact version and shown on the abstract page https://arxiv.org/abs/2605.20695v1. Full licence notice: \texttt{licenses/arxiv-2605.20695-CC-BY-4.0.txt}.\par\smallskip

\noindent\textit{Editorial scope.} Counted unit: the article's lemma L:unit\_expansion (source0.tex lines 427--438). The article's complete original proof of it is delivered with it (source0.tex lines 587--617, one \texttt{proof} environment of 31 lines, which the article places away from the statement). No mathematics is written, repaired or re-derived: the statement and the proof are the article's own lines, in the article's own notation, and no retained line is substituted or re-flowed.  No further source-local context is needed: every cross-reference of every retained line resolves inside the two delivered spans.  Nothing was omitted from any retained span, so the package carries no omission record.  No retained line contains a citation, so the wrapper carries no bibliography and no bibliography entry.  No retained line contains a figure environment, an image-inclusion command or drawing code, so no image is delivered and the package declares no asset dependency.  Editorial formatting only, in addition: an AMS \texttt{amsart} preamble that restates the article's own declarations and macros used by the retained lines, namely the environments \texttt{lemma} on separate counters, together with the standard packages those lines need.  The retained environments print in this document as Lemma 1 in order of appearance, whereas the article prints the same items with the numbering of its own full document.  The complete native source of the article as distributed in the arXiv e-print for this exact version is bundled unchanged as \texttt{sources/201010/source0.tex}.
\begin{lemma}\label{L:unit_expansion}
    Let $f$ be a positive integer, and $0<\delta\leq 1$.  Let $\Lambda$ be a full rank lattice in $\C^f$, such that for all non-zero $x\in \Lambda$, at least one coordinate $x_i$ of $x$ has $|x_i|\geq \delta$.  Further assume that the projection of $\Lambda$ onto one of the coordinates of $\C^f$ is injective.  
    Let $v\geq\delta^{-2}\operatorname{covol}(\Lambda)^{1/f}$.  
    
    Suppose that $|U_\Lambda|\geq u^f$ for some $u>0$.
    Then for every $R\geq 2$,
    there exists a translate $a+\Lambda$ of $\Lambda$ such that
    the set $(a+\Lambda)\cap B_R$ projected onto a coordinate gives a point set $\mc{P}$ in the plane with 
    $$
     2\nu(\mc{P})\geq \left(\frac{u \pi R^2}{4v\delta^{2}}\right)^f \textrm{ and } |\mc{P}|\leq\left(\frac{9R^2}{\delta^2}\right)^{f}.
    $$
\end{lemma}

\begin{proof}[Proof of Lemma~\ref{L:unit_expansion}]
By averaging over $a\in \C^f/\Lambda$, there exists a choice of $a$ such that
\begin{equation*}
|(a+\Lambda)\cap B_{R-1}|
\ge \left(\frac{\pi (R-1)^2}{\covol(\Lambda)^{1/f}}\right)^f.
\end{equation*}
Fix an injective coordinate projection $\mathcal{P}$
of the bounded window $W := (a+\Lambda)\cap B_R$.
The translation-by-$U_\Lambda$ argument
described before Lemma~\ref{L:unit_expansion} implies that
\begin{equation*}
2\nu(\mathcal{P}) \ge |U_\Lambda||(a+\Lambda)\cap B_{R-1}|
\ge u^f \left(\frac{\pi (R-1)^2}{\covol(\Lambda)^{1/f}}\right)^f
\ge \left(\frac{u \pi R^2}{4v\delta^{2}}\right)^f,
\end{equation*}
since $R-1 \ge R/2$.
% Note: This means (R-1)^2 \ge R^2/4.
On the other hand, $\Lambda$ is $\delta$-separated
in the sup-norm on $\C^f$,
so the sumset (Minkowski sum) of $W$ and $B_{\delta/2}$
has volume $\ge |W| |B_{\delta/2}|$,
and $\le |B_{R+\delta/2}| \le |B_{3R/2}|$.
Thus
\begin{equation*}
|\mathcal{P}| = |W| \le \frac{|B_{3R/2}|}{|B_{\delta/2}|}
= \left(\frac{9R^2}{\delta^2}\right)^{f}. \qedhere
\end{equation*}
%(It might also be possible to give a lower bound on
%$|\mathcal{P}| = |W| = |(a+\Lambda)\cap B_R|$,
%but we do not need to do so.)
\end{proof}

\end{document}
